\documentclass[10pt,twocolumn]{article}
\usepackage[margin=0.85in]{geometry}
\usepackage{amsmath,amssymb,amsthm}
\usepackage{fancyhdr}
\usepackage{booktabs}
\usepackage{array}
\usepackage{tabularx}
\usepackage{hyperref}
\usepackage{enumitem}
\usepackage{xcolor}

\title{\textbf{Causal Quantum Gravity:\\
A Graph-Laplacian Spine for Unified Quantum--Relativistic\\
Dynamics and a Discrete Causal Signature}}
\author{Yaoharee Lahtee\\
\small Open Civil Science Initiative}
\date{July 5, 2026}

\pagestyle{fancy}
\fancyhf{}
\fancyhead[L]{\textit{\footnotesize Causal Quantum Gravity}}
\fancyhead[R]{\textit{\footnotesize preprint -- not peer reviewed}}
\fancyfoot[C]{\thepage}
\renewcommand{\headrulewidth}{0.4pt}

\newcolumntype{Y}{>{\raggedright\arraybackslash}X}

\begin{document}
\twocolumn[
  \begin{@twocolumnfalse}
    \maketitle
    \begin{abstract}
    \noindent We report a tier-honest, machine-checked (Coq) result within the
    \emph{readout-not-truth} discrete-substrate framework: a single graph-Laplacian
    ``mother equation'' $M\partial^2\Phi + D\partial\Phi + K L_R\Phi + \nabla V(\Phi) = J-\eta$
    yields, without importing external physics formulas, (i) a quantum energy-dispersion
    relation identical in algebraic form to the relativistic (Klein--Gordon-family)
    mass-shell condition, and (ii) a discrete causal/Lorentzian bilinear form whose
    continuum limit is the d'Alembertian $\Box$. We prove, for the first time in this
    program, that these two readouts are not independent: the quantum dispersion
    condition is \emph{exactly} (pure $\mathbb{Q}$-algebra, no continuum limit) the
    vanishing condition of the same boost-invariant wave operator that governs the
    relativistic branch --- one equation, two readouts. We further report an
    exhaustive, adversarially-audited finding that general relativity (continuum
    curvature, the Schwarzschild solution, Einstein's field equations) is \emph{not}
    derivable from this root, and argue --- consistent with the framework's own
    injected-infinity diagnostic --- that this is the philosophically correct outcome,
    not a gap. In its place we exhibit a genuinely native, axiom-free discrete
    curvature object (Forman--Ricci curvature on the graph) and prove its exact
    algebraic link to an independently-proven stability theorem. A companion
    finite-domain Perfectly-Matched-Layer numerical bridge reproduces a literature
    quasinormal-mode frequency to $\sim$0.1--1.2\%, honestly tiered as
    \texttt{finite\_diagnostic}, not \texttt{Th\_coqc}. Responding to an independent
    adversarial referee review by proving more rather than retreating any claim, we
    further report six results promoting previously informal stances to theorems: an
    exact frequency/UV ceiling forced by the graph's own maximum degree, an exact
    ``no local creation'' energy-balance law, a Schr\"odinger-shaped skew-adjoint
    first-order skeleton, a causal sign-construction theorem (with an honestly
    disclosed partial closure), and a discrete Noether theorem; and one further
    result showing that growing the graph itself is a native, non-continuum discrete
    analog of cosmological expansion --- curving space negatively, conserving energy
    exactly, and diluting spectral modes by a closed-form law, while explicitly not
    claiming to resolve any general-relativistic cosmological discrepancy. A further
    six files add an action-stationarity identity, a curvature-energy-horizon
    structure, and a four-step spectral-convergence ladder resolving the
    flat/diagonal-metric case of the mother equation's continuum limit in
    arbitrary dimension. A further eleven files add a discrete disk-before-lock
    analog, close a seed-to-spectral-ceiling chain, bridge to an algebraic
    track (curvature-Noether invariance, a rotation-matrix identity, a
    golden-ratio-adjacent pentagon spectrum), a graph-native boundary-energy fact and a
    degree-curvature bound, a general tensor-reconstruction fact with a
    curvature bridge conditional on one stated hypothesis, and a momentum-
    optimizer/kernel-energy identification. We further compose several of these
    theorems --- none individually about mass --- into a curvature-limited
    mass bound conditional on a single disclosed, non-theorem identification,
    whose contrapositive states that sufficiently large mass forces negative
    combinatorial curvature; a companion synthesis note included as a
    supplement gives the full derivation and literature positioning. Two
    further files restate the degree/frequency-ceiling family in
    entropy-flavored language (explicitly disclosed as a naming choice, not a
    physical claim) and prove a graph-native boundary-screening theorem
    (information confined to a region's strict interior cannot reach a fixed
    exterior weighting), each explicitly scoped against holographic or
    general-relativistic over-reading. All told:
    69 Coq files, 439 theorem certificates (428 axiom-free \texttt{Th\_coqc}, 11 \texttt{+reals}),
    every one machine-checked. Every claim below
    is accompanied by its tier, its exact source theorem, and a reproduction
    command.
    \end{abstract}

    \textbf{Keywords:} discrete quantum gravity; graph Laplacian; retained-readout
    framework; Klein--Gordon unification; Forman--Ricci curvature; discrete Noether
    theorem; graph growth; Coq formal verification; quasinormal modes; reproducibility.

    \textit{\footnotesize Preprint notice. This manuscript is a preprint (v1.0) and
    has not undergone formal peer review. It has, however, undergone one internal
    editorial pass and one independent adversarial technical review (\S\ref{sec:attack}),
    both logged. Definitions, proofs, and numerical results are reproducible from the
    accompanying repository and remain provisional until independently reproduced by
    a third party.}
    \vspace{0.5em}
    \hrule
    \vspace{1em}
  \end{@twocolumnfalse}
]

\section{Introduction}
This manuscript reports work within a larger research program (the ``readout-not-truth''
framework) whose central commitment is that a formal system should never assert more
than it has machine-checked: a claim is \texttt{Th\_coqc} only if \texttt{coqc} accepts
it and \texttt{Print Assumptions} reports the global context closed over $\mathbb{Q}$;
\texttt{+reals} if it additionally depends on Coq's disclosed Reals axioms;
\texttt{finite\_diagnostic} if it is a reproducible numerical measurement, not a proof;
\texttt{Dr} if it is an interpretive stance; \texttt{Open} if it is an admitted gap.
Collapsing these tiers --- describing a measurement as a proof, or a stance as a
derivation --- is treated as the framework's cardinal error.

The object under study is a single \emph{mother equation} on a finite graph
(Eq.~\ref{eq:mother}), built from one primitive, $\delta_R$ (retained difference), via
a graph Laplacian $L_R$. This paper asks, and answers honestly: what, exactly, follows
from this one equation, and what must be imported?

\section{Claim Tiers}
\label{sec:tiers}
\begin{table}[h]
\centering
\small
\begin{tabularx}{\linewidth}{@{}lY@{}}
\toprule
\textbf{Tier} & \textbf{Meaning} \\
\midrule
\texttt{Th\_coqc} & Machine-checked, axiom-free over $\mathbb{Q}$; \texttt{Print Assumptions} reports the context closed. \\
\texttt{+reals} & \texttt{Th\_coqc} but depends on Coq's disclosed Reals axioms (\texttt{sig\_forall\_dec}, \texttt{functional\_extensionality\_dep}). \\
\texttt{finite\_diagnostic} & A reproducible numerical measurement; not a proof. \\
\texttt{Dr} & An interpretive stance, not machine-checked. \\
\texttt{Open} & An admitted, undischarged gap. \\
\bottomrule
\end{tabularx}
\caption{Claim tiers used throughout this manuscript.}
\label{tab:tiers}
\end{table}

\section{The Mother Equation}
\label{sec:mother}
\begin{equation}
M\partial^2\Phi + D\partial\Phi + K L_R\Phi + \nabla V(\Phi) = J - \eta
\label{eq:mother}
\end{equation}
where $L_R$ is a weighted graph Laplacian built from edge data
$(u_e, v_e, w_e)$, $M,D,K$ are inertia/dissipation/coupling scalars, and $J-\eta$ is a
drive-minus-loss source term. All results below concern the \emph{conservative}
sector ($D=0$, $\nabla V=0$, $J=\eta$) unless stated otherwise.

\textbf{Where $L_R$'s shape and $D$ come from, together.} \texttt{InfoAsymmetricSeedTrifurcation.v}
(\texttt{Th\_coqc}, axiom-free) shows $L_R$ and $D$ need not be specified independently. An
asymmetric directed seed $R_0$ decomposes \emph{exactly} into a diagonal part, a symmetric
off-diagonal part, and an antisymmetric off-diagonal part; the antisymmetric part's sign is
forced by \texttt{nat}'s own trichotomy, not posited. Imposing on the \emph{whole} seed the same
two axioms that already force $L_R$'s own shape (zero row-sum, non-positive off-diagonal) --
rather than only on its symmetric part -- forces
$D_i = \deg_i(W) - \lambda\cdot\mathrm{circ}_i(\mathrm{ord})$: the symmetric-coupling degree at
node $i$ minus the seed's own net directed circulation, down to two free primitives ($W\geq0$,
$\lambda$) instead of the three ($W$, $\lambda$, $D$) a reader might otherwise assume Eq.~\ref{eq:mother}
needs. This identity, and the row-sum axiom it rests on, hold not only for the fixed 3-node witness
used in the exploratory arc but for an \emph{arbitrary} finite vertex set
(\texttt{InfoSeedArbitraryNForcing.v}, \texttt{Th\_coqc}) -- proved once by induction on an arbitrary
\texttt{NoDup} vertex list via one reusable lemma, not re-checked size by size. A companion result
(\texttt{InfoSeedCirculationArbitraryN.v}, \texttt{Th\_coqc}) closes the remaining gap: the seed's net
circulation $\sum_i\mathrm{circ}_i(\mathrm{ord})$ also sums to exactly zero for an arbitrary vertex
set, via a discrete Fubini lemma (swapping the outer and inner sum of a double sum, with the summand's
two arguments swapped to match, leaves the total unchanged) applied to the antisymmetric part directly
-- for any antisymmetric function on any finite vertex set, the double sum over all ordered pairs
equals its own negation, hence is exactly zero. \emph{This is a statement about how many vertices $R_0$
is built on, not about the spatial dimension of Eq.~\ref{eq:mother}'s Laplacian}; it is unrelated to the
$n$-dimensional convergence question of \S\ref{sec:nd}. $M$ and $K$ remain exactly as free as before;
this result concerns only $L_R$'s shape and $D$.

\section{Branch 1: Quantum Readout}
\label{sec:quantum}
\textbf{Source.} \texttt{formal/InfoSchrodinger.v} (extracted from the audited core;
\texttt{Th\_coqc}). A temporal mode gives $\partial^2\to-\omega^2$; on an
$L_R$-eigenmode of eigenvalue $\lambda$ the spine residual is
\begin{equation}
K\lambda - M\omega^2 = 0 \iff M\omega^2 = K\lambda.
\label{eq:dispersion}
\end{equation}
Composing with the Planck--Einstein relation $E=\hbar\omega$ \cite{planck1900,einstein1905a} gives
\begin{equation}
E^2 M = \hbar^2 K\lambda.
\label{eq:energy}
\end{equation}
\textbf{Validation (\texttt{finite\_diagnostic}).} Feeding the $C_6$ (benzene-ring)
graph Laplacian spectrum through Eq.~\ref{eq:energy} reproduces the qualitative
structure of H\"uckel molecular-orbital theory \cite{huckel1931} (adjacency spectrum
$\{2,1,1,-1,-1,-2\}$, degeneracy pattern $1$-$2$-$2$-$1$); the resonance energy
$2\beta$ is graph-theoretic, its numeric value in eV requires the empirical
$\beta$. This is a consistency demonstration, not an independent confirmation of a
fitted constant.

\section{Branch 2: Special-Relativistic Readout}
\label{sec:sr}
\textbf{Source.} \texttt{formal/InfoLorentz.v}, \texttt{InfoLorentzContinuum.v}
(\texttt{Th\_coqc} Tier-0 and \texttt{+reals} respectively; both independently
re-verified via \texttt{Print Assumptions} in this work, though the underlying
proofs were authored 2026-06-27, prior to this manuscript's other results).
A signed bilinear form on graph edges,
\begin{equation}
\mathrm{cform}_{\mathrm{sgn}}(x,y) = \sum_e \mathrm{sgn}(e)\, w_e\, \delta_x(e)\,\delta_y(e),
\label{eq:causalform}
\end{equation}
is self-adjoint, reduces to $L_R$'s own quadratic form when $\mathrm{sgn}\equiv+1$
(Eq.~\ref{eq:euclred}), and is invariant under permutation of the edge list
(Eq.~\ref{eq:perminv}):
\begin{align}
\mathrm{cform}_{+1}(x,y) &= \mathrm{info\_form}(x,y), \label{eq:euclred}\\
\mathrm{cform}_{\mathrm{sgn}}(x,y;\pi(E)) &= \mathrm{cform}_{\mathrm{sgn}}(x,y;E). \label{eq:perminv}
\end{align}
The continuum limit of the corresponding signed second-difference operator is proved
(\texttt{+reals}) to be the d'Alembertian
\begin{equation}
\Box = -\partial_{tt} + \partial_{xx}.
\label{eq:box}
\end{equation}
\textbf{Scope.} The sign function $\mathrm{sgn}$ in Eq.~\ref{eq:causalform} is
universally quantified in the source theorems, not itself derived from a partial
causal order on the graph nodes; a genuine derivation of an indefinite signature from
$\delta_R$'s own (currently linear) order relation is an open research direction
(\S\ref{sec:open}), not claimed here. Separately, the specific Lorentz boost formula
$\gamma(t-vx)$ \cite{lorentz1904,einstein1905b} is an imported \texttt{Definition},
verified (not derived) to leave $\Box$ invariant.

\section{Unification: Branches 1 and 2 Are One Equation}
\label{sec:unification}
This section reports the manuscript's principal new result. Define, following the
already-proven boost-invariant exact quadratic wave operator
$\mathrm{box\_quad}(a_{tt},a_{xx}) := -2a_{tt}+2a_{xx}$, and the spine residual
$R(M,K,\omega^2,\lambda) := K\lambda-M\omega^2$ from Eq.~\ref{eq:dispersion}. Under
the exact reparametrization $a_{tt}:=M\omega^2/2$, $a_{xx}:=K\lambda/2$:
\begin{equation}
\mathrm{box\_quad}\!\left(\tfrac{M\omega^2}{2},\tfrac{K\lambda}{2}\right) = R(M,K,\omega^2,\lambda).
\label{eq:unify1}
\end{equation}
This is a pure ring identity (\texttt{ring} tactic, zero-line proof); it is not a
coincidence of notation but a statement that both objects were built from the same
signed second-difference structure. Consequently,
\begin{equation}
M\omega^2 = K\lambda \iff \mathrm{box\_quad}\!\left(\tfrac{M\omega^2}{2},\tfrac{K\lambda}{2}\right) = 0,
\label{eq:unify2}
\end{equation}
i.e.\ \textbf{the quantum dispersion condition (Branch 1) is exactly the vanishing of
the same wave operator (Branch 2)}. Because $\mathrm{box\_quad}$ is already proven
boost-invariant, this condition is preserved under the identical Lorentz boost:
\begin{equation}
\mathrm{box\_quad}\big(c_{tt}(g,v,\ldots),\,c_{xx}(g,v,\ldots)\big) = 0
\label{eq:unify3}
\end{equation}
whenever Eq.~\ref{eq:dispersion} holds and $g^2(1-v^2)=1$. All three equations are
\texttt{Th\_coqc} (\texttt{formal/InfoQuantumRelativityUnification.v}).

\textbf{Interpretation.} The dispersion relation $E^2M=\hbar^2K\lambda$
(Eq.~\ref{eq:energy}) is quadratic in $E$, structurally matching the relativistic
mass-shell relation $E^2=(pc)^2+(mc^2)^2$, not the linear-in-$E$ non-relativistic
Schr\"odinger dispersion. Equations~\ref{eq:unify1}--\ref{eq:unify3} show this is not
an incidental resemblance: it is because the mother equation's own second-order-in-time
structure is, algebraically, the same object as the relativistic wave operator whose
boost invariance Branch 2 already establishes. Branches 1 and 2 are therefore one
readout of one equation, not two separately-argued branches sharing a name.

\section{Branch 3: General Relativity --- Honestly Not Derived}
\label{sec:gr}
An exhaustive audit of every gravity-touching module in the source repository found
no path from the mother equation to continuum general relativity; every module
either imports the Schwarzschild metric factor, Einstein's field equations, the Unruh
temperature, or the Bekenstein--Hawking entropy as an external \texttt{Definition}
(self-disclosed in each module's own header), or formalizes generic tensor algebra
over free variables never connected to $L_R$.

Eight independent attempts to close this gap were tested this program and refuted or
left open: a bare definitional alias; a dimensionally-forced numerological coincidence
(shown to equally ``confirm'' an unrelated quantity, hence non-discriminating); an
axiom-restatement; a structurally mass-independent prediction contradicted by a
ten-billion-fold mass comparison against the true $\kappa\propto1/M$ scaling; a partial
WKB match; two hyperboloidal-compactification attempts that hit a genuine
\emph{irregular} singular point at spatial infinity (an essential singularity,
re-injecting exactly the actual-infinity artifact this framework's own diagnostic
refuses); and one success (\S\ref{sec:qnm}) that shares only a discretization
\emph{method}, not an equation, with the mother equation.

\textbf{Conclusion.} Continuum general relativity --- curvature defined via
derivative limits on a complete manifold --- is, by this framework's own
injected-infinity diagnostic, a non-readout. Its absence from this derivation is
therefore the philosophically correct outcome, matching the unresolved state of every
other discrete-substrate research program surveyed (causal sets \cite{bombelli1987},
the Wolfram Physics Project \cite{wolfram2020}, Regge calculus \cite{regge1961}).

\section{Branch 4: Native Discrete Curvature}
\label{sec:curvature}
\textbf{Source.} \texttt{formal/InfoDiscreteGraphCurvature.v} (\texttt{Th\_coqc},
axiom-free). Rather than chasing continuum curvature, we exhibit a discrete curvature
notion --- Forman--Ricci curvature \cite{forman2003} --- that requires no derivative,
limit, or manifold, and is already a native function of $L_R$'s own edge/degree data:
\begin{equation}
F(u,v) = 4 - \deg(u) - \deg(v).
\label{eq:forman}
\end{equation}
We prove that under uniform edge weight $w$, the weighted degree
$\mathrm{wdeg}$ (from an independently-proven coercivity/stability theorem) equals
$w\cdot\deg$ exactly:
\begin{equation}
\mathrm{wdeg}(i) = w\cdot\deg(i) \implies C_{\mathrm{safe}}V_{\max}\,\mathrm{wdeg}(i) = C_{\mathrm{safe}}V_{\max}w\deg(i).
\label{eq:curvlink}
\end{equation}
The same degree count that sets an edge's curvature (Eq.~\ref{eq:forman}) also sets,
by exact substitution, the dissipation a node requires for the spine to remain
coercive. We note, following adversarial review, that this link is exact but shallow:
any monotone function of degree would exhibit the same shared dependence
(\S\ref{sec:attack}).

\textbf{The same curvature, now tied to the seed of \S\ref{sec:mother}.}
\texttt{InfoSeedCurvatureIntegration.v} (\texttt{Th\_coqc}) builds Eq.~\ref{eq:forman}'s
curvature from a uniform-weight instantiation of the \S\ref{sec:mother} seed, rather than a
separately-posited weighted graph. Every node of this instantiation has degree $2$, so every
edge is Forman-flat ($F=0$) -- an honest, unglamorous consequence of a $2$-regular graph, not a
numerically striking one, following the same pattern already noted for the $C_6$/benzene case
elsewhere in this paper. The weighted degree $\mathrm{wdeg}$ is computed, not assumed, to equal
exactly the degree term appearing in the seed's forced-$D$ identity of \S\ref{sec:mother} -- a
check on two independently-defined quantities, not a coincidence built in by construction. As
with Eq.~\ref{eq:curvlink}, no underlying theorem becomes tighter or more general; what is new is
this concrete tie between curvature and the forced-$D$ construction, on the same seed.

\section{Bridge: Quasinormal-Mode Numerics}
\label{sec:qnm}
\textbf{Source.} \texttt{formal/InfoQuantumGravityRootBridge.v} (\texttt{+reals}) +
\texttt{scripts/verify\_quantum\_gravity\_root\_bridge.py}
(\texttt{finite\_diagnostic}). The Regge--Wheeler equation \cite{regge1957}
\begin{equation}
\frac{d^2\psi}{dr_*^2} + [\omega^2 - V(r)]\psi = 0
\label{eq:rw}
\end{equation}
built on the already-verified Schwarzschild metric factor $f(r)=1-2M/r$, is
discretized as a finite path-graph Laplacian eigenvalue problem with a
Perfectly-Matched-Layer absorbing boundary \cite{berenger1994} --- no point at
infinity anywhere. This converges ($N=1600$--$6400$) to
\begin{equation}
M\omega \approx 0.4841 - 0.0956\,i,
\label{eq:qnmresult}
\end{equation}
matching the literature scalar $\ell=2$, $n=0$ fundamental quasinormal-mode value
$M\omega\approx0.4836-0.0968\,i$ \cite{leaver1985,berti2009} to $\sim$0.1\% (real
part) / $\sim$1.2\% (imaginary part), robust across grid resolution, absorbing-layer
strength, and domain size.

\section{A Strengthening Campaign: Six Results Closing Referee-Flagged Gaps}
\label{sec:campaign}
Following the same referee review, six further results were added, each promoting a
previously informal (\texttt{Dr}-tier) stance to a \texttt{Th\_coqc} theorem --- by
proving more, not retreating any existing claim. Every file: pre-verified by
exact-rational numerical testing before authoring, then independently compiled and
\texttt{Print Assumptions}-checked here; all Tier-0 axiom-free; zero \texttt{funext},
zero classical axioms, zero \texttt{admit}.

\textbf{Frequency ceiling and CFL window (\texttt{InfoSpectralCeiling.v},
\texttt{InfoRecurrenceEnergy.v}, \texttt{InfoQuantumFrequencyCeiling.v}, 20 theorems).}
A pure degree-sum argument --- no eigenvalue theory, no square root --- gives
$\lambda \le 2 d_{\max}$, a Rayleigh-quotient form of the Gershgorin bound. Composed
with Branch 1's dispersion relation, this gives a hard frequency ceiling
\begin{equation}
M\omega^2 \le K \cdot (2 d_{\max}),
\label{eq:freqceiling}
\end{equation}
forced by the graph's own maximum degree, not asserted as a physical constant.
Composed further with an exact discrete-leapfrog Lyapunov identity, the same degree
bound guarantees both a CFL-style stability window $0\le a\le4$ and energy-non-
increasing dynamics under dissipation. This replaces an earlier interpretive
``$\tau_c$ floor'' stance with an exact structural theorem about what sets it.

\textbf{No local creation (\texttt{InfoGraphFluxBalance.v}, 8 theorems).} A discrete
divergence theorem plus Green's identity (summation-by-parts) prove that $L_R$'s
coupling term contributes exactly zero to the mother equation's total energy budget
--- it telescopes to zero across every edge. Energy can only change via dissipation
($\le0$) or source work.

\textbf{A Schr\"odinger-shaped skeleton (\texttt{InfoCompanionSkew.v}, 5 theorems).}
Writing the conservative sector as a first-order companion system $\Psi=(x,v)$ with
generator $B(x,v)=(Mv,-KLx)$, the generator is proved skew-adjoint under the energy
inner product $\langle(x,v),(y,w)\rangle_E = K\cdot\mathrm{gform}(x,y)+M\langle
v,w\rangle$, moving every state exactly orthogonal to its own energy level set. This
is the algebraic core of norm-preserving first-order evolution --- the $i\hbar
\partial_t\psi=H\psi$ skeleton without $i$, without $\mathbb{C}$, without $\sqrt{\ }$.
(One of the five, \texttt{green\_identity}, is a helper lemma re-proved inside this
file so it checks standalone, not a fifth headline result beyond the four above.)
This "without $i$" skew-adjoint skeleton is a separate, weaker claim than
\S\ref{sec:algebra}'s later, narrower observation that the mother equation's
own stepper matrix literally IS the standard $i$-rotation matrix at one
specific stability-window parameter value: that later claim never invokes
$\mathbb{C}$ or Hilbert-space structure either, and is scoped to a single
real $2\times2$ matrix identity, not to the Schr\"odinger-shaped skeleton
above. The two do not conflict, but should not be read as reinforcing each
other without this note.

\textbf{Causal signature (\texttt{InfoCausalSignature.v}, 7 theorems).} \S\ref{sec:sr} disclosed
that $\mathrm{sgn}$ is a free parameter. This file proves a sign function can instead
be \emph{constructed} from a relation's comparability (comparable pairs $\to -1$,
incomparable pairs $\to +1$), yielding an exact PSD-minus-PSD split, two lightcone-
style cone inequalities, and, on a concrete star-graph example, a genuine rational-
congruence witness of Minkowski type $(1,3)$. \textbf{This does not fully close the
gap.} This repository's own causal order (built from a strict total order on its
domain, proved order-isomorphic to $\mathbb{N}$) admits no incomparable pairs at all
--- applying the construction to it would degenerate to an all-timelike signature,
not the indefinite structure demonstrated on the constructed example. The
construction is real and general; it does not yet connect to this repository's own
specific causal order, which would require a genuinely richer (non-total) structure.

\textbf{Discrete Noether (\texttt{InfoGraphNoether.v}, 7 theorems).} Given a graph
automorphism $\sigma$, the Laplacian commutes with it, every quadratic structure
built from $L_R$ is $\sigma$-invariant, and the antisymmetric pairing
$W(p,q)=\sum_i p(\sigma i)q_i - p_i q(\sigma i)$ is exactly conserved under the
conservative step: genuine Noether shape, the inertial term cancels pointwise and the
coupling term dies by equivariance plus summation-by-parts. A concrete instance (the
6-cycle with a rotation automorphism) confirms the theorem is non-vacuous.

\section{Graph Growth: A Native Discrete Expansion Analog}
\label{sec:growth}
A natural question, raised independently, is whether this framework can address a
real cosmological discrepancy such as the Hubble tension (the $\sim5$--$6\sigma$
disagreement between early-universe and local-distance-ladder measurements of
$H_0$). The honest answer is no: that is a general-relativistic/FRW-cosmology
problem, and \S\ref{sec:gr} establishes this repository derives no general-relativity
content at all. Forcing a connection would repeat the same category error as the
eight refuted GR attempts of \S\ref{sec:gr}.

Following the same discipline that produced Branch 4 (curvature, \S\ref{sec:curvature}
--- a native discrete object in place of an unreachable continuum target), a natural
discrete analog of \emph{expansion} was investigated instead: growing the graph
itself. \texttt{InfoGraphGrowth.v} (13 theorems, axiom-free) proves, for the mother
equation's own graph structure:

\textbf{Growth curves space negatively.} Adding an edge $e^*$ to a graph strictly
decreases the Forman curvature (\S\ref{sec:curvature}) of every edge adjacent to
$e^*$ by at least~1, exactly:
\begin{equation}
F_{e^*::E}(e) = F_E(e) - \mathrm{share}(e^*,u) - \mathrm{share}(e^*,v).
\label{eq:growthcurve}
\end{equation}
``Space creation curves space negatively'' is a combinatorial identity, not an
analogy.

\textbf{Growth has an exact energy account.} Adding an edge changes the vector energy
balance (\S\ref{sec:campaign}) by exactly $h^2K$ times the product of the field
difference across the new edge; a field continuous across the new edge changes the
energy not at all (adiabatic growth).

\textbf{The order-collapse obstruction, proved.} \S\ref{sec:campaign}'s causal-
signature gap --- this repository's causal order is total, so it cannot produce an
indefinite signature --- is here promoted from an observation to a theorem: for
\emph{any} graph where every edge is comparable under the order, the signed form
collapses to $-\mathrm{gdiag} \le 0$ everywhere. ``Space is incomparability; a total
order has no space'' is now exact, not descriptive. The qualitative statement ---
that a total causal order admits no spatial structure --- is established folklore
in causal set theory \cite{surya2019}; a literature check found no prior published
proof of this specific curvature-collapse form, but the underlying intuition is
not new to this work and is cited as such.

\textbf{An exact dilution law.} On a ramp/path growth sequence, the Rayleigh quotient
of a specific test state has closed form $12/(n(n+1))$ --- division-free throughout,
via an $\mathbb{N}\to\mathbb{Q}$ power-sum injection. Paired with the frequency
ceiling (Eq.~\ref{eq:freqceiling}), the growing graph brackets its own spectrum from
both sides: a ceiling set by degree, a floor that dilutes as $\sim 1/n^2$.

\textbf{Scope (\texttt{finite\_diagnostic}, honestly bounded).} No exponential or de
Sitter-type dilution law was found; only power laws were observed, and a star-growth
sequence does not dilute at all --- the dilution law is a property of the growth
\emph{topology}, not of growth per se. This is a native, non-continuum, discrete
expansion analog: growth simultaneously curves space negatively, has exact energy
accounting, and dilutes spectral modes with a closed-form law. It is not a
cosmological expansion rate, does not reproduce $H_0$ under any interpretation, and
Hubble-parameter matching remains exactly the non-readout target this repository
does not chase.

\section{Action Stationarity and a Curvature-Energy-Horizon Structure}
\label{sec:action}
Two further results (\texttt{InfoActionStationarity.v}, 8 theorems, and
\texttt{InfoCurvatureBalance.v}, 12 theorems, both Tier-0 axiom-free) are reported
here plainly, for exactly what they prove --- neither inflated nor understated.

\textbf{The mother equation is an exact stationarity readout.} Define the discrete
window action $S=\sum_n\left[\frac{M}{2}\left(\frac{\Delta x_n}{h}\right)^2 -
\frac{K}{2}\mathrm{gform}(x_n,x_n)\right]$. Perturbing a single node by $s$ changes
the action by an exact quadratic identity (\texttt{action\_polarization}); on the
conservative step, the linear coefficient of that perturbation --- the discrete
analogue of $\partial S/\partial x_n$, computed via polarization in $\mathbb{Q}$,
with no limit --- vanishes identically at every node
(\texttt{mother\_step\_stationary}). The mother equation \emph{is} the discrete
Euler--Lagrange stationarity condition of this action, exactly, not by analogy.
Two further facts sharpen this: \texttt{damped\_step\_balance} shows dissipation is
\emph{exactly} the obstruction to stationarity (the first variation under the
damped step is $2sc(q_i-p_i)$, vanishing only when $c=0$); \texttt{strict\_descent}
shows that whenever the step residual is nonzero, an explicit variation
($s=R/C$) strictly decreases the action --- non-solutions are never selected, a
genuine, exact argmin content, not a metaphor.

\textbf{What this does and does not establish.} This is a real, native fact about
the mother equation, worth stating in full: an argmin/stationarity reading of the
dynamics that was previously informal is now an exact theorem. It is reported here
at full strength. What it does \emph{not} do is narrow any gap toward general
relativity specifically. Being expressible as the stationarity of an action is a
property shared by essentially every Lagrangian physical theory --- classical
mechanics, Maxwell's equations, and general relativity itself all share it. Two
theories both admitting a variational formulation is not evidence that one derives
the other; it is the most generic structural property a physical law can have.
This result stands on its own merits and is reported at that value, no more and no
less. A literature check (2026-07-05) found this file's specific combination ---
one theorem for the matter/field side on a fixed graph, a separate theorem for an
edge-addition cost-benefit law on the graph's own topology --- is not identical to
any single prior construction, but is closely related to Trugenberger's
Combinatorial Quantum Gravity \cite{trugenberger2017}, which varies a statistical
action over dynamical random graphs where matter arises as curvature excitations
and geometry emerges from graph condensation, with the Einstein--Hilbert action
recovered in a geometric phase. The relation is noted explicitly rather than left
for a reader to discover: this file's two separately-proved stationarity
statements are a different formalization choice from Trugenberger's single
combined action, not an independent invention of the underlying idea of varying
both matter and graph topology together.

\textbf{A curvature-energy-horizon structure.} \texttt{InfoDiscreteGraphCurvature.v}'s
degree identity composes exactly with the growth results of \S\ref{sec:growth}:
$\deg(u)+\deg(v) = 4-F(u,v)$ (\texttt{edge\_deg\_curvature}), so the coercivity
dissipation threshold from \S\ref{sec:curvature} is an exact affine function of
Forman curvature (\texttt{horizon\_threshold\_curvature}) --- more negative
curvature forces a strictly higher threshold (\texttt{threshold\_antitone}). Taking
the per-edge retention benefit to be affine in curvature, $b(e)=\alpha+\beta F(e)$
--- an explicit modeling choice, not independently derived from anything else in
this framework --- gives an exact two-way iff between edge-retention and curvature
(\texttt{curvature\_balance\_add}, \texttt{curvature\_balance\_rem}): under that
choice, energy strain and curvature determine each other exactly. Separately, and
without depending on that ansatz, a genuine dichotomy is proved directly from the
dissipation-threshold structure: \texttt{no\_escape}/\texttt{no\_escape\_strict}
show that once per-node dissipation exceeds every admissible control's budget, no
choice can prevent energy from being non-increasing; \texttt{repair\_exists} shows
the converse --- below that threshold, a control exists that keeps a local mode
from decaying. This gives a principled (non-circular) answer to an earlier open
question about how to fix the per-node horizon threshold: it is exactly the point
at which dissipation exceeds every admissible selection's budget, and this budget
is set by curvature. \texttt{vacuum\_joint\_stationary} confirms non-vacuity: a
constant field is a joint stationary point of both the field and (under the
curvature-affine ansatz) the geometry side, on every graph.

\textbf{Scope, stated plainly.} The horizon dichotomy (\texttt{no\_escape}/
\texttt{repair\_exists}) is unconditional and exact. The curvature-balance iff
(\texttt{curvature\_balance\_add/rem}) is exact conditional on the affine-benefit
ansatz --- a real theorem about a chosen model, not a derivation of why that model
is the correct one. Neither result derives Einstein's field equations, the
specific coupling constant, or dynamical geometry sourced by matter in general;
neither is claimed to.

\section{The $n$-Dimensional Spectral Convergence Ladder}
\label{sec:nd}
Whether $L_R$ spectrally converges to a continuum Laplace--Beltrami operator in
more than one dimension has been an open frontier since Branch 4
(\S\ref{sec:curvature}). Four further files (\texttt{InfoProductSpectrum.v}, 7
theorems, Tier-0; \texttt{InfoContinuumLimit\_nD.v}, 4 theorems, Tier-1;
\texttt{InfoWeightedReadout.v}, 2 theorems, Tier-1; \texttt{InfoCrossTermDominance.v},
9 theorems, Tier-0) factor this into four steps rather than attempting it as one
result, each tagged at its true tier.

\textbf{Step 1 (Tier-0): Kronecker/Rayleigh additivity.} An $n$-dimensional grid is
the Cartesian product of $1$-D chains, and a product graph's quadratic form
separates exactly: $\mathrm{pgform}(x\otimes y) =
\mathrm{gform}_G(x)\|y\|^2+\|x\|^2\mathrm{gform}_H(y)$
(\texttt{form\_separates}), with the analogous exact identity for the norm
(\texttt{norm\_separates}). Consequently a Rayleigh pair on each factor lifts to a
Rayleigh pair on the product with eigenvalues adding (\texttt{rayleigh\_additive})
--- pure $\mathbb{Q}$ combinatorics, no eigenvalue theory or continuum limit
invoked.

\textbf{Step 2 (Tier-1, \texttt{sig\_forall\_dec}+\texttt{funext} only, confirmed
no \texttt{classic} leak): the flat multi-axis readout.} Reusing this repository's
own $1$-D symmetric-second-difference machinery (\texttt{tends0}, verbatim from
the audited core) with two new closure lemmas (additivity, scaling),
\texttt{weighted\_two\_axis\_readout} shows a weighted sum of per-axis
second-difference quotients converges to the same weighted sum of the per-axis
limits. \texttt{box\_2d\_readout} (weights $-1,1$) reproduces this repository's
own \texttt{lorentz\_box} construction as an instance --- confirming the extension
is consistent with the existing kernel, not a divergent reformulation.
\texttt{laplacian\_2d\_readout}/\texttt{laplacian\_3d\_readout} give the flat
$n$-D Laplacian.

\textbf{Step 3 (Tier-1, same axiom profile): the variable-coefficient case.} The
flux-form stencil $a(x{+}h/2)(f(x{+}h){-}f(x)) - a(x{-}h/2)(f(x){-}f(x{-}h))$,
divided by $h^2$, is shown (\texttt{weighted\_second\_difference\_limit}) to
converge to $a(x)f''(x)+a'(x)f'(x) = (af')'(x)$ --- the divergence-form operator
emerges as a readout of the discrete flux stencil, not assumed.
\texttt{constant\_coefficient\_instance} confirms consistency with Step 2.

\textbf{Step 4 (Tier-0): the off-diagonal dominance criterion, sharp both ways.}
An anisotropic (cross-term-bearing) quadratic form is representable as a
nonnegative-weighted graph form --- i.e.\ as a genuine graph Laplacian --- if and
only if it is diagonally dominant (\texttt{repr\_pos\_iff\_dominant},
\texttt{repr\_neg\_iff\_dominant}). Sufficiency is a direct decomposition;
necessity is extracted by evaluating any hypothetical representation at three
probe vectors, forcing the cross-edge weight to a specific value that goes
negative exactly when dominance fails --- a genuine obstruction, not an artifact
of one decomposition choice.

\textbf{What this closes and what remains open, by design.} Steps 1--4 together
give the diagonal-metric Laplace--Beltrami operator class
$\sum_i\partial_i(a_i(x)\partial_i u)$ --- covering Schwarzschild-radial, FRW, and
essentially every metric weak-field tests use --- with a sharp criterion for when
off-diagonal terms admit a graph representation at all. What remains open, stated
at its narrowest: spectral convergence for a fully general (non-dominant)
anisotropic metric, and general-metric convergence on a genuine manifold. The
latter is a known, substantial result in the literature
(Burago--Ivanov--Kurylev, 2014) \cite{bik2014}; it is cited here as prior
literature, not mechanized, and is not claimed as this work's contribution.

\section{A Discrete Disk-Before-Lock Analog}
\label{sec:disklock}
A discrete, graph-native analog of an unrelated preprint's continuum claim
("Disk Before Mass," a PDE result stated with a provenance map but no worked
derivation, hence \texttt{Dr} tier and not relied upon here) is built from
scratch using only this repository's own retention machinery. Under an
explicit \emph{anisotropy hypothesis} --- edges classified ``perpendicular''
or ``parallel,'' with the perpendicular class's benefit always strictly below
its strain, and the parallel class's always at or above --- a perpendicular
candidate edge is shown to strictly worsen the retention functional while a
parallel candidate never does, from any current graph state, unconditionally.
This is an immediate, per-decision fact, not an asymptotic claim about many
growth steps: \texttt{InfoDiskBeforeLock.v} (5~theorems, \texttt{Th\_coqc}).
The anisotropy hypothesis itself is disclosed as input data, not derived.

\section{Closing the Seed-to-Ceiling Chain}
\label{sec:seedchain}
A separate card (Supplement) tracks a chain from the seed's own retention
criterion through curvature to an internal spectral ceiling. Two remaining
gaps in that chain are closed here. \texttt{InfoGrowthFold.v}
(5~theorems, \texttt{Th\_coqc}) generalizes the single-edge curvature-shift
law of \S\ref{sec:campaign} to an arbitrary \emph{list} of added edges via an
exact, order-independent closed form (\texttt{forman\_fold}), and shows
retention only ever bends curvature in one direction regardless of growth
order. \texttt{InfoCeilingMonotone.v} (6~theorems, \texttt{Th\_coqc}) shows
growth never decreases the quadratic form a Rayleigh quotient is built from,
that every exact Rayleigh level achieved before growth is still dominated
after growth, and that the degree data feeding the spectral ceiling of
\S\ref{sec:campaign} only grows --- explicitly disclosing what is \emph{not}
claimed (individual eigenvalue monotonicity/interlacing, which needs
spectral theory not invoked here). Together these make the seed-through-
ceiling portion of that chain gap-free up to one disclosed ansatz and one
pre-existing, direction-undecidable axiom relating an internal oscillation
period to mass --- neither of which is discharged by this section, and
neither of which is claimed to be.

\textbf{A third gap closed: the seed of \S\ref{sec:mother} feeds an existing spectral
ceiling.} \texttt{InfoSeedTauRelFloor.v} (\texttt{Th\_coqc}) applies the Anderson--Morley
bound \cite{andersonmorley1985} $\lambda\le 4-F_{\min}$ (machine-checked elsewhere in this repository) to a concrete
eigenpair of the \S\ref{sec:mother} seed's own unweighted carrier: eigenvector $(1,-1,0)$ is
verified, not assumed, to satisfy the node identity with eigenvalue $3$. The genuinely new
content is a reciprocal-monotonicity lemma turning the curvature ceiling on $\lambda$ into a
genuine \emph{lower bound} on a relaxation time $1/\lambda$ -- the actual falsifiable content
(relaxation faster than this floor would violate the cited bound itself). Stated openly: not
numerically dramatic on this small, flat carrier (the floor is unsaturated, $1/4$ against an
actual value of $1/3$); this closes a third gap in the chain described above, not a
strengthening of the ceiling itself.

\section{Algebraic Bridges: Curvature Symmetry, Rotation, and a Pentagon Spectrum}
\label{sec:algebra}
Three results connect the curvature/growth track to a separately-developed
algebraic track (complex roots of unity, the golden ratio) via concrete
shared objects, not analogy. \texttt{InfoCurvatureNoether.v}
(3~theorems, \texttt{Th\_coqc}) proves Forman curvature is invariant under
any graph automorphism with a correspondingly permuted edge list, extending
an existing Noether-style invariance result (proved there only for the
Laplacian-built quadratic form) to curvature specifically for the first
time; a concrete instance reuses an existing order-6 rotation witness, with
the honest caveat that this particular witness graph is regular, so its
curvature is trivially constant there --- the general theorem is what
matters for non-regular graphs. \texttt{InfoModeRotation.v}
(8~theorems, \texttt{Th\_coqc}) shows the mother equation's own leapfrog/CFL
stepper matrix, at the specific stability-window parameter value where its
trace is zero, is \emph{literally} the standard matrix representation of
multiplication by $i$ (verified: the same squares-to-negation, period-4
identity already known for $i$), and that two further window points give
periods 6 and 3 --- instances of the crystallographic restriction theorem,
cited as prior literature for the general claim (only $n\in\{1,2,3,4,6\}$
are reachable by a rational stepper parameter), not re-derived here.
\texttt{InfoPentagonSpectrum.v} (5~theorems, \texttt{Th\_coqc}) proves the
5-cycle graph Laplacian satisfies an exact cubic operator identity whose
nonzero-eigenvalue factor has the same defining-relation shape as the golden
ratio's own characterizing equation (both are minimal polynomials of
elements of $\mathbb{Q}(\sqrt5)$), and exhibits all four distinct
eigenvalues of the 6-cycle (the same graph used as the automorphism witness
above) via explicit rational eigenvectors --- the 6-cycle is, in this
precise sense, the last cycle whose full spectrum stays rational; the
5-cycle is the first whose spectrum leaves $\mathbb{Q}$.

\textbf{A fourth bridge: the seed's antisymmetry IS discrete torsion, literally.}
\texttt{InfoSeedTorsionIsSkewOff.v} (6~theorems, \texttt{Th\_coqc}) specializes the standard
differential-geometry torsion definition (the antisymmetric part of a connection's
coefficients) to the rank-1/scalar-bundle case, where the \S\ref{sec:mother} seed $R_0$ itself
IS the connection matrix: torsion equals exactly twice the seed's own antisymmetric part, a
genuine instantiation of the textbook definition, not an analogy. A symmetric connection is
torsion-free; for the forced seed of \S\ref{sec:mother}, torsion is exactly the directional
scalar $\lambda$ times the sign forced by \texttt{nat}'s trichotomy -- the seed's own
directionality \emph{is}, literally, its torsion. \texttt{InfoSeedTorsionGroupAndRankN.v}
(5~theorems, \texttt{Th\_coqc}) checks a coboundary connection (already proven elsewhere in this
repository to carry genuine holonomy curvature) is separately group-torsion-free -- holonomy
curvature and torsion are two independent kinds of connection non-triviality, confirmed rather
than assumed -- and extends the rank-1 identity to $n$ independent decoupled seeds sharing one
$\lambda$, explicitly flagged as the decoupled special case. \texttt{InfoSeedTorsionGenuineMixing.v}
(4~theorems, \texttt{Th\_coqc}) closes the harder general case: a genuinely mixed connection
$\Gamma(k,i,j) := T(k,i)R(i,j)$, where output component $k$ and base pair $(i,j)$ entangle
through a shared table $T$, not $n$ independent copies. A concrete non-separable $T$ (the
Kronecker delta) gives a computed witness where the zero/nonzero pattern of the resulting torsion
itself changes with $k$; the key theorem rigorously proves no single scalar reduces one $k$'s
torsion to another's -- a genuine impossibility result, not an illustration.

\section{A Graph-Native Boundary-Energy Fact and a Degree-Curvature Bound}
\label{sec:arealaw}
Two small, self-contained results, explicitly not claimed as more than what
they state. \texttt{InfoAreaLaw.v} (3~theorems, \texttt{Th\_coqc}) proves
that if a field is constant across a region of a graph, every edge strictly
interior to that region contributes exactly zero to the quadratic energy
form --- all of a uniform region's energy comes from its boundary. This
shares a qualitative flavor with the Bekenstein--Hawking holographic bound
(``boundary determines the budget, not the interior'') but is not that
bound and makes no claim about entropy, metric area, or general relativity;
sharing a qualitative flavor is not sharing a theorem.
\texttt{InfoDegreeFromCurvature.v} (2~theorems, \texttt{Th\_coqc}) proves a
node's degree is bounded by $4$ minus the local Forman curvature of any
incident edge, and, given a curvature floor over a graph, that the maximum
degree is bounded by $4$ minus that floor --- a purely combinatorial fact
usable to feed the spectral ceiling of \S\ref{sec:campaign} with a
curvature-derived value. Stated plainly: given the file's own definition
\texttt{forman(e) := 4 - deg(u) - deg(v)}, the local bound is an immediate
one-line rearrangement (equivalent to non-negativity of degree); the
contribution is packaging this as a directly-usable input to the spectral
ceiling, not a deep new fact about curvature.

\section{General Tensor Reconstruction and a Conditional Curvature Bridge}
\label{sec:tensorframe}
\texttt{InfoTensorFrame.v} (7~theorems, \texttt{Th\_coqc}, division-free
over $\mathbb{Q}$) proves a general, dimension-independent linear-algebra
fact: a symmetric bilinear form on an $n$-dimensional space is fully and
uniquely reconstructible from its $n$ axis (diagonal) evaluations plus its
$\binom{n}{2}$ face-diagonal evaluations, via the polarization identity ---
the counting identity $2n+n(n-1)=n(n+1)$ confirms no slack in any
dimension, and in $4$D, $4$ axes $+$ $6$ face-diagonals $=10$ matches a
rank-2 symmetric tensor's component count exactly. This is stated plainly
as an \emph{abstract} linear-algebra fact, deliberately not identified with
general relativity's metric or stress-energy tensor: no metric signature,
no coordinate transformation law, and no field equation is claimed or used
here, and the fact holds for a symmetric bilinear form on \emph{any}
$n$-dimensional $\mathbb{Q}$-vector space, gravitational or not. It also
does not by itself establish that any graph actually used elsewhere in this
work contains face-diagonal edges (the $n$-D ladder of \S\ref{sec:nd} uses
axis-aligned edges only).

\texttt{InfoStrainTensorBridge.v} (11~theorems, \texttt{Th\_coqc}) connects
this to curvature on a stated condition: on a $2{\times}2$ graph cell with a
diagonal candidate edge, the strain on that edge equals the tensor-frame
diagonal evaluation \emph{exactly when the cell is affine} (no discrete
Hessian cross term), and differs from it by an exact, fully disclosed
defect term in general. Under that affine condition, the seed's own
retention criterion for the diagonal candidate is shown to be, verbatim, a
price on a tensor-frame evaluation (\texttt{diagonal\_retention\_prices\_
tensor}) --- curvature and this tensor-reconstruction fact are connected by
a real theorem for the affine case, not merely both being expressible on
the same graph. This is \emph{not} an unconditional closure; for a
non-affine cell the disclosed defect term must be carried. The same file
separately proves that any graph's quadratic form splits exactly into
inside, outside, and cut contributions against any region
(\texttt{gform\_screen\_partition}, unconditional), and, under one further
disclosed lens (benefit per retained distinction taken proportional to a
temperature-like parameter), that the retention criterion becomes a
discrete Clausius-type inequality, both directions exact. None of this
constitutes or is claimed to constitute Jacobson's thermodynamic
derivation of the Einstein equation \cite{jacobson1995}; it is a small,
honestly-tiered fragment sharing a qualitative shape with one step of that
program, explicitly not the program itself.

\section{Momentum Optimizer Theory Coincides With the Kernel's Energy Theory on a Quadratic Mode}
\label{sec:optimizer}
\texttt{InfoOptimizerWindow.v} (10~theorems, \texttt{Th\_coqc}, division-free
over $\mathbb{Q}$) shows the classical heavy-ball/momentum iteration on a
quadratic mode is, under an exact reparametrization, the same dissipative
two-step recurrence already proved for the mother equation's own leapfrog
stepper in \S\ref{sec:campaign}. The reparametrized Lyapunov functional
decrements by exactly one nonnegative square per step; a momentum
dichotomy (never increases below a threshold, exactly conserved at the
threshold, never decreases above it) is proved as a theorem; and an
a-priori bound on every iterate is obtained from initial data alone, with
no spectral computation or run-time tuning, inside the classical stability
window. Divergence outside that window, convergence rates, and anything
about non-quadratic objectives are explicitly not claimed.

\section{A Composed Reading: Curvature-Limited Mass}
\label{sec:masschain}
A companion synthesis note (\texttt{paper/mass\_note.tex} in this repository,
also independently authored by the present author, cross-checked against this
repository at the exact commit this manuscript itself was built from) composes
several of the theorems above --- none of them individually about mass --- into
a mass chain: retention $\to$ degree $\to$ curvature $\to$ an internal-mode
frequency ceiling $\to$ mass, through exactly one disclosed, non-theorem
identification. We state the composition here at main-text weight because it
is a genuine, checkable consequence of results already established above, not
a new file; the companion note carries the full derivation, literature
positioning, and a 21-file SHA-256-pinned artifact table.

The Anderson--Morley eigenvalue ceiling \cite{andersonmorley1985} is now
mechanized directly, in witness form over $\mathbb{Q}$
(\texttt{InfoSpectralCeilingSharp.v}, 2~theorems plus 4~supporting lemmas,
all \texttt{Th\_coqc}; \texttt{anderson\_morley\_witness}):
for any exact eigenpair, some edge $(u,v)$ satisfies
$\lvert\lambda\rvert\le\deg_E(u)+\deg_E(v)$, proved by taking the edge
maximizing $\lvert x_u-x_v\rvert$ and comparing the eigen-equation at its two
endpoints directly --- no edge-space Gershgorin matrix, no division. Composed
with the exact identity $\deg_E(u)+\deg_E(v)=4-F(e)$
(\texttt{InfoDegreeFromCurvature.v}) under a curvature floor, this gives the
\emph{sharp} curvature ceiling directly
(\texttt{sharp\_curvature\_ceiling}), $\lambda\le4-F_{\min}$ --- a factor
$\sqrt2$ tighter, at the frequency level, than the coarser $2d_{\max}$ route
this manuscript mechanized earlier (\S\ref{sec:campaign},
\texttt{InfoSpectralCeiling.v}), which remains correct but is no longer the
sharpest available bound. The mass ceiling below now uses the sharp form
throughout:
\begin{equation}
\omega_{\mathrm{int}}^2 \;\le\; \frac{K}{M}\,(4-F_{\min}),
\label{eq:massomega}
\end{equation}
a direct algebraic substitution of two already-\texttt{Th\_coqc} facts, itself
\texttt{Th\_coqc} (no additional Coq file needed: both inputs are theorems,
and the composition is one inequality chain). Given exactly one imported,
disclosed identification --- $\tau_c:=1/\omega_{\mathrm{int}}$ together with
the programme's own mass--clock postulate $m=\hbar/(2\tau_c c^2)$ (a
programme postulate, not a published result, tagged \texttt{[AX/import]} in
the companion note, selected over two explicitly rejected sign-analysis
alternatives there) --- Eq.~\ref{eq:massomega} becomes a curvature-limited
mass bound,
\begin{equation}
m \;\le\; \frac{\hbar}{2c^2}\sqrt{\frac{K}{M}\,(4-F_{\min})}.
\label{eq:massceilmain}
\end{equation}
This bound is \texttt{Th\_coqc} \emph{conditional on} the one imported
identification, and holds as stated \emph{for any excitation that is an exact
mode of the graph carrying it}: the inequality it rearranges is fully
mechanized, but Eq.~\ref{eq:massceilmain} itself is not a standalone Coq
theorem (no \texttt{Print Assumptions} block asserts it directly) --- it is
the companion note's own arithmetic on top of \eqref{eq:massomega}, disclosed
as such there. A second, independent caveat applies to the physically
motivating case of an \emph{approximately localized} excitation restricted to
a support subgraph rather than an exact mode of the whole graph: the
companion note states plainly that this restriction is, at present, only a
numerical finding (a diagnostic mode carrying 99\% of its weight on a dense
cluster across an ensemble of random graphs), \texttt{finite\_diagnostic}, not
a kernel theorem even granting the ansatz above --- and we repeat that
distinction here rather than let the theorem-conditional-on-ansatz framing
imply more than it does for the localized case, which is exactly the case the
``curvature license'' reading below is really about.
Its contrapositive is the note's sharper claim: mass exceeding the
$F_{\min}=0$ cap $(\hbar/2c^2)\sqrt{4K/M}$ forces negative curvature
somewhere on the excitation's support --- curvature is not sourced by mass
here, it is a capacity mass requires, the reverse of general relativity's
habitual direction of explanation, and inherits both caveats above. Both the
bound and its contrapositive additionally inherit a disclosed dependence on
whether Forman curvature deserves the name at all --- the same open question
already carried by this manuscript's Open Problem~3 (\S\ref{sec:open}) and
named \textsc{ob-forman-ricci} in the companion note --- neither claim is
stronger than that dependence allows, and the companion note states this
explicitly rather than eliding it.

\section{An Entropy-Language Restatement, and a Boundary-Screening Theorem}
\label{sec:entropyscreen}
Two further small, self-contained results, each explicitly scoped to state
no more than what it proves.

\texttt{InfoEntropyLicense.v} (6~theorems, \texttt{Th\_coqc}, division-free
over $\mathbb{Q}$) restates the degree/frequency-ceiling family of
\S\ref{sec:campaign} in entropy-flavored language: defining a per-node
quantity $S_{\mathrm{loc}}(s_0,E,i):=s_0\cdot\deg_E(i)$ for an arbitrary
weight $s_0$, the handshake identity becomes
$\sum_i S_{\mathrm{loc}}(s_0,E,i)=2s_0\lvert E\rvert$; $S_{\mathrm{loc}}$ is
proved monotone non-decreasing under retention-only graph growth (the same
$E_s{+}{+}E$ shape as \S\ref{sec:growth}); and, given the ``achieving''
hypothesis $2d_{\max}=\deg_E(u)+\deg_E(v)$ (an explicit hypothesis, of
exactly the same kind \S\ref{sec:campaign}'s own spectral ceiling already
takes for $d_{\max}$), the frequency-ceiling inequality $M\omega^2\le
K\cdot2d_{\max}$ composes \emph{genuinely} with the entropy identity into a
bound on the entropy-language quantity itself,
$M\omega^2\le K\cdot(S_{\mathrm{loc}}(s_0,E,u)+S_{\mathrm{loc}}(s_0,E,v))/s_0$
--- not a mere side-by-side pairing of two unrelated facts, but one
inequality substituted into the other. A further rearrangement gives a
lower bound on $S_{\mathrm{loc}}$ at the degree-realizing node in terms of a
plain-$\mathbb{Q}$ energy relation. We state plainly what this is and is
not: $s_0$ is an arbitrary $\mathbb{Q}$-parameter and $S_{\mathrm{loc}}$ is,
by definition, a rescaling of degree; calling it ``entropy'' is a naming
choice in the same family as the qualitative Bekenstein--Hawking
resemblance already disclaimed for the boundary-energy fact of
\S\ref{sec:arealaw}, not a new physical claim, not the Bekenstein
bound~\cite{bekenstein1973}, and not an instance of Wheeler's ``it from
bit''~\cite{wheeler1990}. No thermodynamic, statistical-mechanical, or
information-theoretic (in the Shannon or Szil\'ard--Landauer
bit-erasure~\cite{szilard1929,landauer1961} sense) identification of $s_0$
or $S_{\mathrm{loc}}$ is claimed or used anywhere in the file.

\texttt{InfoBoundaryScreening.v} (5~theorems, \texttt{Th\_coqc}) proves a
graph-native screening fact on top of \texttt{InfoStrainTensorBridge.v}'s
region-partition identity of \S\ref{sec:tensorframe}
(\texttt{gform\_screen\_partition}): calling the outside-plus-cut part of
the quadratic form against a region \emph{Exterior}, and calling a node
\emph{strictly interior} when it is in the region and every incident edge
of the graph stays inside the region, Exterior's value is exactly
unchanged by any change to a field's values confined to non-interior
nodes---two fields that agree everywhere off the strict interior give
\emph{identical} Exterior values, unconditionally and exactly, for any
graph, region, and pair of fields (\texttt{exterior\_invisible\_to\_
interior\_swap}). The proof is a direct case analysis on which of the
region-partition's three edge classes (inside/outside/cut) each edge falls
into: every edge contributing to Exterior has at least one endpoint outside
the region, and that very edge already witnesses that both its endpoints
fail the strict-interior test, so the fixed weighting is untouched when the
field is redefined there. We state plainly what this does and does not
establish. It \emph{does} show that, on this graph's own native quadratic
form, information about the field's values in the strict interior cannot
propagate into Exterior --- a purely combinatorial, unconditional fact
about which edges carry a given field to the outside. It does \emph{not}
establish anything about physical entropy, the holographic
principle~\cite{thooft1993,susskind1995}, a Bekenstein-Hawking area
law~\cite{bekenstein1973}, or general relativity; ``boundary'',
``screening'', ``Exterior'', and ``capacity'' (the trivially field-independent
cut-edge count, stated for completeness with a one-line proof) are
graph-native names chosen for readability, not physical identifications, in
the same spirit as \S\ref{sec:arealaw}'s explicit disclaimer that a
qualitative resemblance to the Bekenstein--Hawking bound is not that bound.

\section{Opening the Nonlinear Sector, and a Readability-Boundary Instance}
\label{sec:cubicreadability}
Two further small results, one opening genuinely new ground for this
kernel, one closing one more instance of an existing algebraic-restriction
story.

\texttt{InfoCubicLinearization.v} (6~theorems, \texttt{Th\_coqc},
division-free over $\mathbb{Q}$) is the first file in this arc to give the
mother equation's $\nabla V(\Phi)$ slot a concrete, nonzero, mechanized
instance: every prior result in this manuscript works in the conservative,
purely quadratic sector ($\nabla V=0$). Adding a single on-site cubic term
$-g\Phi^3$ (from $V(\Phi)=(g/4)\Phi^4$) to the existing linear force, the
file proves the exact one-step polarization identity for the force
difference under a field perturbation, extracts the linear-order
coefficient shift $q(i):=3g\,\Phi(i)^2$ that the perturbation acquires from
the background, and its sign ($q\ge0$ for $g\ge0$, strictly positive on the
background's support for $g>0$). What is \emph{not} claimed, stated in the
file's own header: any time-averaging or homogenization step (the passage
from a pointwise coefficient shift to an effective-medium mass reading),
and any propagation or spectral consequence --- those remain explicitly
disclosed, not mechanized. A companion numerical probe (not part of this
manuscript's own claim set, reported in the project's supplementary
material) found that locking energy into a zero-group-velocity mode of a
region measurably slows a passing low-frequency packet only when $g\ne0$
(a linear, $g=0$ control shows no effect, consistent with superposition);
this is offered as motivation for why the $\nabla V$ term matters, not as
something this file itself establishes.

\texttt{InfoReadabilityBoundary.v} (3~theorems, \texttt{Th\_coqc}, plus one
worked example and one small composing remark) proves that
no rational parameter of \S\ref{sec:algebra}'s stepper matrix makes its
5th iterate the identity: the relevant polynomial factors exactly as
$-(a^2{-}5a{+}5)(a^2{-}3a{+}1)$ (an exact ring identity), and both
quadratic factors have discriminant~5, so a rational root would force a
rational square root of 5 --- proved impossible from scratch via integer
descent. This is one further instance in the same family as
\S\ref{sec:algebra}'s periods 3, 4, and 6, but is stated narrowly: it does
\emph{not} establish that $\{1,2,3,4,6\}$ is the complete set of rational
periods reachable by this stepper, which would require the general
crystallographic restriction theorem (cited as literature, not mechanized
here) --- an earlier, broader draft of this result was reviewed, found not
to follow from the rest, and dropped rather than weakened into an
inaccurate statement.

\section{Reviewer Attack Response}
\label{sec:attack}
This section pre-empts questions an adversarial referee will ask, following an
independent technical review of an earlier draft.

\noindent\textit{Is $E^2M=\hbar^2K\lambda$ really quantum mechanics?} No, not
non-relativistic Schr\"odinger mechanics (correctly flagged: that theory is
first-order in $E$). It is a relativistic (Klein--Gordon-family) dispersion, and
\S\ref{sec:unification} proves this is exact, not a relabeling.

\noindent\textit{Is the H\"uckel match a real validation?} It is a consistency
demonstration of graph-topology sensitivity (confirmed via a non-aromatic control),
not an independent confirmation of a fitted physical constant (\S\ref{sec:quantum}).

\noindent\textit{Is Eq.~\ref{eq:perminv} really ``frame covariance''?} It is
invariance under edge-list relabeling, a bookkeeping symmetry, not physical boost
covariance; the latter (Eq.~\ref{eq:box}'s invariance under the imported boost
formula) is stated and proved separately and is not conflated with it here.

\noindent\textit{Is $\mathrm{sgn}$ in Eq.~\ref{eq:causalform} actually derived from
the causal order $\prec$?} No --- it is a free parameter in the source theorems. This
manuscript states that explicitly (\S\ref{sec:sr}) rather than implying otherwise, and
lists deriving it as open work (\S\ref{sec:open}).

\noindent\textit{Is the curvature--dissipation link (Eq.~\ref{eq:curvlink})
non-trivial?} It is an exact algebraic identity, but shallow: it holds because both
quantities are stated in terms of degree, not because curvature causes a dissipation
requirement. This is disclosed, not obscured.

\noindent\textit{Are all citations independently checkable?} Every citation in
\S\ref{sec:refs} was checked against standard physics/mathematics references; two
citations flagged as unresolved in an internal draft (an anonymous arXiv preprint and
a mismatched preprint server) have been removed from this version pending independent
verification, rather than retained on trust.

\noindent\textit{Does the frequency ceiling (Eq.~\ref{eq:freqceiling}) claim a
physical Planck-scale minimal length?} No. It is a structural/discretization bound: a
property of $L_R$'s finite maximum degree. It states \emph{that} a ceiling exists and
\emph{exactly} what sets it ($d_{\max}$, $M$, $K$), not a numerical value for any
particular physical system.

\noindent\textit{Is the causal-signature closure (\S\ref{sec:campaign}) actually
complete?} No, and the manuscript says so directly: the construction is proved for an
arbitrary relation with genuine incomparable pairs, demonstrated on a constructed
example, but this repository's own causal order is a total order and cannot host one.
This is reported as a partial result with a named obstruction, not a closed gap.

\noindent\textit{Does the graph-growth result (\S\ref{sec:growth}) address the Hubble
tension or any cosmological discrepancy?} No, explicitly. That is a general-
relativistic problem and \S\ref{sec:gr} establishes no general relativity is derived
here. The growth result is offered only as a native discrete analog of expansion,
with an honestly bounded scope (power-law dilution only; no de Sitter behavior found;
topology-dependent, not universal).

\noindent\textit{Does the boundary-energy fact (\S\ref{sec:arealaw}) establish a
holographic entropy bound?} No. It shares a qualitative flavor with the
Bekenstein--Hawking bound (boundary determines the budget, not the interior) but
proves a purely combinatorial fact about a quadratic energy form on a graph ---
no entropy, no metric area, and no general-relativistic structure appears in the
statement or the proof.

\noindent\textit{Does \S\ref{sec:tensorframe}'s tensor-reconstruction fact
establish anything about physical tensors (metric, stress-energy), and is the
curvature bridge to it unconditional?} No on both counts. The reconstruction
fact is an abstract linear-algebra identity holding for any symmetric bilinear
form on any $n$-dimensional $\mathbb{Q}$-vector space; it is not applied to, or
claimed to be, general relativity's metric or stress-energy tensor. The
curvature bridge (\texttt{diagonal\_retention\_prices\_tensor}) holds only on
the stated affine-cell hypothesis; for a non-affine cell an explicitly disclosed
defect term must be carried, and this is not an unconditional closure of any
kind.

\noindent\textit{Is the curvature-limited mass bound (\S\ref{sec:masschain})
itself a Coq theorem?} No, and this manuscript says so directly. The
inequality it algebraically rearranges (Eq.~\ref{eq:massomega}) is a direct
composition of two already-mechanized theorems and is itself \texttt{Th\_coqc};
the mass bound (Eq.~\ref{eq:massceilmain}) additionally depends on one
imported, non-theorem identification ($\tau_c:=1/\omega_{\mathrm{int}}$ plus
the imported mass--clock relation), so it is conditional, not unconditional
\texttt{Th\_coqc} --- exactly as labeled. No new Coq file mechanizes
Eq.~\ref{eq:massceilmain} as a single theorem. A second, lower-tier caveat
applies specifically to localized (particle-like) excitations restricted to a
support subgraph rather than an exact whole-graph mode: that restriction is
currently only a numerical (\texttt{finite\_diagnostic}) finding, not a
theorem, even granting the ansatz above --- stated in \S\ref{sec:masschain}
and repeated here because it is precisely the case the mass narrative is
about.

\section{Novelty Audit}
\label{sec:novelty}
The individual mathematics throughout --- Gershgorin/Rayleigh bounds, leapfrog
invariants, discrete Green's identities, Forman curvature, lattice Noether
theorems, variational integrators, Kronecker sums, diagonal-dominance criteria ---
is classical or already in the literature, and this repository's own discrete-
substrate-physics framing sits in a crowded field (causal sets, the Wolfram
Physics Project, Regge calculus, loop quantum gravity). A dedicated literature
check (2026-07-05) on the five specific structural identifications introduced in
\S\ref{sec:action}--\ref{sec:growth} found no close prior art for three of them
(the curvature-affine horizon threshold; the horizon dichotomy proved rather than
stipulated; an epistemic-axiom-to-energy-floor bound, not mechanized here), one
close but distinct prior construction requiring explicit citation (the double
matter/geometry stationarity, \S\ref{sec:action}, cf.\ Trugenberger's
Combinatorial Quantum Gravity \cite{trugenberger2017}), and one instance of
well-established qualitative folklore given a specific proved form here for the
first time as far as this check found (the order-collapse theorem,
\S\ref{sec:growth}, cf.\ causal set theory \cite{surya2019}).

\textbf{Scope limit, stated plainly.} This literature check covers only
\S\ref{sec:action}--\ref{sec:growth}. The eleven files and six sections added
later (\S\ref{sec:disklock}--\ref{sec:optimizer}) have NOT been through an
equivalent dedicated check and their novelty is therefore \emph{unaudited}:
several of their individual ingredients are flagged in-line as classical or
adjacent to known results (the crystallographic restriction theorem cited as
prior literature in \S\ref{sec:algebra}; the Bekenstein--Hawking qualitative
resemblance disclaimed in \S\ref{sec:arealaw}; the polarization identity
underlying \S\ref{sec:tensorframe}), but no systematic search for close prior
art was performed for these sections as a whole. \S\ref{sec:masschain}'s
composed mass reading is a partial exception: the companion synthesis note
that develops it independently positions each constituent equation against
ten literature streams by owner and states its own search protocol for the
narrower claim that no prior discrete-gravity programme possesses a
machine-checked kernel of any size; that audit is the companion note's own
and is not re-performed or independently re-verified here.

The seed-asymmetry additions to \S\ref{sec:mother}, \S\ref{sec:curvature}, \S\ref{sec:seedchain},
and \S\ref{sec:algebra} (the $R_0$ decomposition, its arbitrary-$n$ generalization, the seed-tied
curvature and relaxation-time-floor results, and the torsion identification) are, likewise,
\emph{unaudited} for novelty against the discrete-torsion/discrete-differential-geometry
literature specifically -- no dedicated search was performed for prior graph-native torsion
constructions comparable to \S\ref{sec:algebra}'s fourth bridge. Individual mathematical
ingredients (Forman--Ricci curvature, the Anderson--Morley bound \cite{andersonmorley1985}) are
already cited to their own literature elsewhere in this manuscript; what has not been checked is
whether the specific \emph{combination} -- an asymmetric seed forcing both a coupling operator and
its dissipation together, with its antisymmetric part identified as literal torsion -- has close
prior art. This is stated here rather than left for a reader to discover.

The claimed contribution is therefore narrower than "new physics" and is stated
at that level: a single graph operator carrying genuine (not aliased) derivations
--- not results sharing a vocabulary, but a shared equation --- across dispersion,
causal structure, conservation, symmetry, curvature, and variational structure
simultaneously, machine-checked and axiom-free end to end, with a tier system
that separates theorem from stance at every step. This is a claim about
\emph{method and integration}, not about any single physical prediction; the one
structural prediction native to this framework is the frequency/UV ceiling of
\S\ref{sec:campaign} (Eq.~\ref{eq:freqceiling}), which has not been tested against
any experiment.

\section{Main-Text vs.\ Supplement Boundary}
\label{sec:boundary}
This manuscript is self-contained for its theorem-level claims: every stated
equation carries its exact source file and tier, and \S\ref{sec:repro} gives the
exact commands to independently re-check every one. The manuscript is responsible
for the mother equation, the four branch results (\S\ref{sec:quantum}--
\ref{sec:curvature}), the unification result (\S\ref{sec:unification}), the
quasinormal-mode bridge (\S\ref{sec:qnm}), the strengthening campaign
(\S\ref{sec:campaign}), graph growth (\S\ref{sec:growth}), the action-stationarity
and curvature-horizon structure (\S\ref{sec:action}), the $n$-D spectral ladder
(\S\ref{sec:nd}), the disk-before-lock analog (\S\ref{sec:disklock}), the
seed-to-ceiling closure (\S\ref{sec:seedchain}), the algebraic bridges
(\S\ref{sec:algebra}), the boundary-energy fact and degree-curvature bound
(\S\ref{sec:arealaw}), the tensor-reconstruction and conditional curvature
bridge (\S\ref{sec:tensorframe}), the momentum-optimizer identification
(\S\ref{sec:optimizer}), the composed curvature-limited mass reading
(\S\ref{sec:masschain}), the entropy-language restatement and
boundary-screening theorem (\S\ref{sec:entropyscreen}), the reviewer attack
response (\S\ref{sec:attack}), and the reproduction commands.

The companion synthesis note (\texttt{paper/mass\_note.tex}, compiling to
\texttt{paper/mass\_note.pdf}) is a separate, self-contained preprint by the
same author, included here as a supplement because it is built directly on
this repository's own theorems at this manuscript's own commit; it carries
its own abstract, literature review (ten streams, positioned individually),
circularity audit, and open-problem register, and should be read as a
companion document, not as part of this manuscript's own claim set. Where the
two overlap (\S\ref{sec:masschain} above), this manuscript's statement is the
authoritative one for this repository's own claim set; it is not, on every
point, strictly narrower than the companion note's --- \S\ref{sec:masschain}
now carries both of the companion note's caveats on the mass bound (the
ansatz dependence and the separate localization-to-support-subgraph caveat)
precisely so this manuscript is not less conservative than the note it
summarizes.

The companion supplement (\texttt{SUPPLEMENT.md}) is responsible for material
this manuscript summarizes but does not reproduce in full: the complete
dependency DAG in both diagram and prose form; the full log of all eight
attempts to derive general relativity, including the ones only briefly
mentioned in \S\ref{sec:gr}; the complete novelty audit against prior
discrete-substrate literature (\S\ref{sec:novelty} gives only its conclusion);
and the per-citation verification notes recording how each reference in
\S\ref{sec:refs} was checked, including two citations present in an earlier
internal draft and removed after they could not be independently verified
(\S\ref{sec:attack}). A reader who wants the theorem statements and can
reproduce them needs only this manuscript; a reader who wants the full
argument trail, including dead ends, needs the supplement as well.

\section{Reproducibility Protocol}
\label{sec:repro}
\begin{enumerate}[nosep,leftmargin=1.2em]
\item \texttt{git clone} this repository.
\item \texttt{make verify} --- compiles all 69 \texttt{formal/*.v} files via
\texttt{coqc} in dependency order and runs \texttt{Print Assumptions} on every named
theorem (439 total: 428 \texttt{Th\_coqc}, 11 \texttt{+reals}), printing a tier report
(\texttt{Th\_coqc} / \texttt{+reals} per file).
\item \texttt{python3 scripts/verify\_quantum\_gravity\_root\_bridge.py} ---
reproduces the quasinormal-mode convergence table (Eq.~\ref{eq:qnmresult}).
\item Expected result: all Coq files compile with exit code 0; the tier report
matches Table~\ref{tab:tiers}; the QNM script's $N=3200$ row matches
Eq.~\ref{eq:qnmresult} to the printed precision ($N=1600$ is already close,
$0.4838-0.0958i$, but the digits quoted in Eq.~\ref{eq:qnmresult} are the
$N=3200$ row).
\end{enumerate}

\section{Open Problems}
\label{sec:open}
(1) Construct a genuinely non-total (partial, multi-dimensional) causal order on this
repository's graphs, so that the sign-construction theorem of \S\ref{sec:campaign}
produces a non-degenerate indefinite signature on the repository's \emph{own} causal
structure, not only on a constructed example.
(2) Extend the quasinormal-mode bridge (\S\ref{sec:qnm}) to spin-2 perturbations.
(3) Determine whether Ollivier--Ricci curvature \cite{ollivier2009} offers a
non-shallow alternative to Eq.~\ref{eq:curvlink}'s exact-but-shallow link.
(4) Does the dilution law of \S\ref{sec:growth} admit a principled selection rule for
which growth topologies are physically relevant, or is topology-dependence the final
answer?
(5) Does \S\ref{sec:campaign}'s Noether pairing admit a principled dissipative
correction, or is exact conservation strictly a conservative-sector fact?
(6) Discharge, or find a counterexample to, the pre-existing,
direction-undecidable axiom relating an internal oscillation period to mass
that \S\ref{sec:seedchain} leaves untouched --- the seed-to-ceiling chain is
gap-free only up to this axiom.
(7) Derive, rather than disclose as input data, the anisotropy hypothesis of
\S\ref{sec:disklock} (which edges count as ``perpendicular'' vs.\ ``parallel''
and why) from a more primitive graph property, so the disk-before-lock result
no longer depends on an externally-supplied classification.
(8) Prove or refute an inertia cost law (cost of re-addressing the internal
loop of \S\ref{sec:masschain} scales as $\propto m$ under acceleration); the
companion synthesis note states this is currently unproven, not merely
uncited, and flags it as its own furthest interpretive claim.

\section{Conclusion}
One graph-Laplacian root equation genuinely derives, and genuinely unifies, a
quantum dispersion relation and a special-relativistic wave operator. It honestly
does not derive general relativity, for a stated philosophical reason rather than a
numerical shortfall, and offers in its place a native discrete curvature readout
with an exact, if shallow, link to an independently-proven stability result.
Responding to adversarial review, six further theorems promote a frequency ceiling,
an energy-balance law, a first-order unitary-like skeleton, a causal sign
construction (partially closed, honestly), and a discrete Noether law from informal
stances to machine-checked facts; a seventh shows that graph growth is a native,
non-continuum, honestly-bounded analog of expansion, explicitly not a claim about
any physical cosmological discrepancy. Two further results prove the mother
equation is an exact stationarity readout on both its field and its geometry, and
compose the graph's own curvature into an exact horizon threshold and dichotomy ---
reported at full strength for what they establish, and explicitly not as evidence
of narrowing any gap toward general relativity's field equations, since
variational-structure-sharing is generic to essentially all physical theories. A
final four-step ladder closes the flat and diagonal-metric cases of this
framework's $n$-dimensional continuum limit, tagging each step at its true tier and
narrowing the remaining general-metric case to a single, named, literature-cited
gap. Eleven further files extend the reach without ever advancing a claim past
its evidence: a discrete anisotropy result immediate at every decision point,
not asymptotic; two files closing a seed-to-spectral-ceiling chain up to one
disclosed ansatz and one pre-existing, direction-undecidable axiom, neither
discharged here; three files bridging curvature and growth to a
separately-built algebraic track via concrete shared objects (an
automorphism-invariance extension, a literal identity between the mother
equation's own stepper matrix and the standard $i$-rotation matrix, and a
pentagon-graph spectrum landing in the same field as the golden ratio); a
graph-native boundary-energy fact and a degree-curvature bound, each explicitly scoped
against the physical results (holographic entropy, general relativity) they
share only a qualitative flavor with, not a theorem; a general tensor-
reconstruction fact and a curvature bridge to it that is conditional on one
named hypothesis, not unconditional; and a momentum-optimizer identification
inheriting the mother equation's own energy theory by exact reparametrization.
Every claim above carries its tier and its exact reproduction command.

\section*{Acknowledgements}
This work was developed through a human--AI co-thinking workflow, including one
round of independent adversarial technical review external to the author's own
assessment. Any remaining errors are the author's own.

\section*{Author Declarations}
\textbf{Funding.} No specific grant from any funding agency in the public,
commercial, or not-for-profit sectors was received for this research.\\
\textbf{Competing Interests.} The author declares no competing interests.\\
\textbf{AI Assistance Disclosure.} Multiple AI assistants were used, in an
assistant role only, for formal-proof drafting, literature cross-checking, numerical
implementation, and editorial structuring, including independent adversarial review
passes. No individual model is named as a contributor: no single model produced this
work, and the research direction, core ideas, substantive claims, tier assignments,
and final decisions are, and remain, the human author's.\\
\textbf{Data and Code Availability.} All Coq sources, the Python reproduction
script, and this manuscript's source are available in the accompanying repository.

\clearpage
\begin{thebibliography}{99}
\small
\bibitem{planck1900} M. Planck, \"Uber das Gesetz der Energieverteilung im
Normalspectrum, \textit{Annalen der Physik}, 1900.
\bibitem{einstein1905a} A. Einstein, \"Uber einen die Erzeugung und Verwandlung des
Lichtes betreffenden heuristischen Gesichtspunkt, \textit{Annalen der Physik}, 1905.
\bibitem{einstein1905b} A. Einstein, Zur Elektrodynamik bewegter K\"orper,
\textit{Annalen der Physik}, 1905.
\bibitem{lorentz1904} H. A. Lorentz, Electromagnetic phenomena in a system moving
with any velocity smaller than that of light, \textit{Proc. Acad. Science
Amsterdam}, 1904.
\bibitem{huckel1931} E. H\"uckel, Quantentheoretische Beitr\"age zum
Benzolproblem, \textit{Zeitschrift f\"ur Physik}, 1931.
\bibitem{forman2003} R. Forman, Bochner's method for cell complexes and
combinatorial Ricci curvature, \textit{Discrete \& Computational Geometry}, 2003.
\bibitem{ollivier2009} Y. Ollivier, Ricci curvature of Markov chains on metric
spaces, \textit{Journal of Functional Analysis}, 2009.
\bibitem{andersonmorley1985} W. N. Anderson, T. D. Morley, Eigenvalues of the
Laplacian of a graph, \textit{Linear and Multilinear Algebra} 18:141--145, 1985.
\bibitem{regge1957} T. Regge, J. A. Wheeler, Stability of a Schwarzschild
singularity, \textit{Phys. Rev.}, 1957.
\bibitem{leaver1985} E. Leaver, An analytic representation for the quasi-normal
modes of Kerr black holes, \textit{Proc. R. Soc. Lond. A}, 1985.
\bibitem{berti2009} E. Berti, V. Cardoso, A. O. Starinets, Quasinormal modes of
black holes and black branes, \textit{Class. Quantum Grav.}, 2009.
\bibitem{berenger1994} J.-P. B\'erenger, A perfectly matched layer for the
absorption of electromagnetic waves, \textit{J. Comput. Phys.}, 1994.
\bibitem{bombelli1987} L. Bombelli, J. Lee, D. Meyer, R. Sorkin, Space-time as a
causal set, \textit{Phys. Rev. Lett.}, 1987.
\bibitem{wolfram2020} S. Wolfram, \textit{A Project to Find the Fundamental Theory
of Physics}, Wolfram Media, 2020.
\bibitem{regge1961} T. Regge, General relativity without coordinates, \textit{Il
Nuovo Cimento}, 1961.
\bibitem{bik2014} D. Burago, S. Ivanov, Y. Kurylev, A graph discretization of the
Laplace-Beltrami operator, \textit{J. Spectral Theory}, 2014.
\bibitem{trugenberger2017} C. A. Trugenberger, Combinatorial quantum gravity:
geometry from random bits, \textit{JHEP} 09:045, 2017.
\bibitem{surya2019} S. Surya, The causal set approach to quantum gravity,
\textit{Living Reviews in Relativity}, 2019.
\bibitem{jacobson1995} T. Jacobson, Thermodynamics of spacetime: the Einstein
equation of state, \textit{Phys. Rev. Lett.} 75, 1260, 1995.
\bibitem{bekenstein1973} J. D. Bekenstein, Black holes and entropy,
\textit{Phys. Rev. D} 7, 2333, 1973.
\bibitem{wheeler1990} J. A. Wheeler, Information, physics, quantum: the search
for links, in \textit{Complexity, Entropy and the Physics of Information},
Addison-Wesley, 1990.
\bibitem{szilard1929} L. Szil\'ard, \"Uber die Entropieverminderung in einem
thermodynamischen System bei Eingriffen intelligenter Wesen, \textit{Zeitschrift
f\"ur Physik} 53, 840, 1929.
\bibitem{landauer1961} R. Landauer, Irreversibility and heat generation in the
computing process, \textit{IBM J. Res. Dev.} 5, 183, 1961.
\bibitem{thooft1993} G. 't Hooft, Dimensional reduction in quantum gravity, in
\textit{Salamfestschrift}, World Scientific, 1993; arXiv:gr-qc/9310026.
\bibitem{susskind1995} L. Susskind, The world as a hologram, \textit{J. Math.
Phys.} 36, 6377, 1995.
\end{thebibliography}
\label{sec:refs}

\end{document}
